\chapter{Estructura del proyecto}
\label{chap:estructura-proyecto}

\section{Proposito de la organizacion}
\label{sec:estructura-proposito}

La estructura del proyecto separa configuracion, dominios de aplicacion, documentacion, recursos estaticos, archivos media y comandos operativos. Esta separacion favorece el mantenimiento porque cada cambio tiene una zona natural: la configuracion general vive en \archivo{config/}, la regla funcional en \archivo{apps/}, los assets compartidos en \archivo{static/} y la documentacion tecnica en \archivo{docs/}.

El criterio dominante es organizar por dominio de negocio. Registro, torneo y comunicaciones no estan mezclados en una sola aplicacion Django. Esto permite revisar un flujo sin recorrer todo el repositorio: si el cambio afecta inscripciones, empieza en \archivo{apps/participants}; si afecta grupos, batallas o podio, empieza en \archivo{apps/tournament}; si afecta plantillas, lotes o correos, empieza en \archivo{apps/invitations}.

\section{Arbol principal}
\label{sec:estructura-arbol}

El siguiente arbol resume la estructura mantenible del proyecto. Se excluyen artefactos locales o generados como \archivo{__pycache__/}, \archivo{.venv/}, \archivo{.env}, bases SQLite, \archivo{media/} operativo y caches.

\begin{lstlisting}[caption={Arbol principal del proyecto},label={lst:estructura-arbol-principal}]
pre_explotaglobos/
  apps/
    common/
    core/
    invitations/
    participants/
    tournament/
  config/
  docs/
    auditoria_tecnica/
    correos/
    manual_tecnico/
    screenshots/
  logos_patrocinadores/
  static/
    css/
    docs/
    js/
    sponsors/
  manage.py
  requirements.txt
  README.md
  .env.example
\end{lstlisting}

El arbol completo puede ampliarse en \cref{anexo:arbol-proyecto}. En este capitulo se prioriza explicar responsabilidades y relaciones para que la estructura sea util durante mantenimiento.

\section{Raiz del proyecto}
\label{sec:estructura-raiz}

\begin{longtable}{p{4cm}p{9cm}}
\toprule
Elemento & Responsabilidad \\
\midrule
\archivo{manage.py} & Entrada de comandos Django. Se usa para validaciones, migraciones, superusuario, servidor local y comandos de gestion. \\
\archivo{requirements.txt} & Dependencias Python declaradas del proyecto. Es la fuente para reproducir el entorno de aplicacion. \\
\archivo{README.md} & Descripcion general existente del proyecto. No sustituye este Manual Tecnico. \\
\archivo{.env.example} & Plantilla segura de variables de entorno con placeholders. Debe copiarse a \archivo{.env} y completarse fuera del control de versiones. \\
\archivo{.gitignore} & Excluye entornos virtuales, secretos, SQLite, media, staticfiles, caches y listados reales de comunicaciones. \\
\bottomrule
\end{longtable}

La raiz conserva pocos archivos porque la logica de aplicacion esta delegada a \archivo{apps/}. Esta convencion reduce el acoplamiento accidental: los archivos de arranque y dependencias no compiten con modelos, vistas o servicios.

\section{Configuracion central}
\label{sec:estructura-config}

La carpeta \archivo{config/} contiene la configuracion central Django.

\begin{longtable}{p{4cm}p{9cm}}
\toprule
Archivo & Responsabilidad \\
\midrule
\archivo{settings.py} & Carga variables de entorno, define apps instaladas, middleware, base de datos, static/media, correo, LibreOffice y controles de seguridad. \\
\archivo{urls.py} & Distribuye rutas raiz hacia admin, sitio publico, registro, comunicaciones y torneo. \\
\archivo{asgi.py} & Punto de entrada ASGI preparado por Django. \\
\archivo{wsgi.py} & Punto de entrada WSGI preparado para servidores compatibles. \\
\bottomrule
\end{longtable}

\archivo{settings.py} es especialmente importante porque aplica reglas de seguridad por entorno. La configuracion bloquea SQLite si \variable{DEBUG=False}, exige \variable{DJANGO_SECRET_KEY} y \variable{DJANGO_ALLOWED_HOSTS} fuera de debug y activa controles como cookies seguras, redireccion SSL y HSTS en entornos no debug.

\section{Aplicaciones locales}
\label{sec:estructura-apps}

Las aplicaciones locales viven bajo \archivo{apps/}. Cada una sigue convenciones Django, pero el proyecto agrega servicios de dominio para evitar que las vistas acumulen regla de negocio compleja.

\subsection{apps.common}

\archivo{apps.common} contiene \clase{TimeStampedModel}, modelo abstracto reutilizable con campos temporales. Su existencia evita repetir definiciones de \variable{created_at} y \variable{updated_at} en modelos de participantes, torneo e invitaciones. No tiene migraciones numeradas confirmadas, solo paquete \archivo{migrations/}.

\subsection{apps.core}

\archivo{apps.core} agrupa la parte publica general:

\begin{itemize}
  \item \archivo{views.py}: vistas \funcion{sponsor_catalog}, \funcion{home}, \funcion{rules_page} y \funcion{sponsors_page}.
  \item \archivo{urls.py}: rutas publicas \ruta{/}, \ruta{/reglas/} y \ruta{/patrocinadores/}.
  \item \archivo{templates/core/}: base publica, home, reglas y patrocinadores.
\end{itemize}

Su responsabilidad es presentar informacion del evento y no ejecutar operaciones internas de torneo. Esta frontera evita que la capa publica termine mezclada con acciones protegidas.

\subsection{apps.participants}

\archivo{apps.participants} contiene registro publico, asistencia y sincronizacion hacia equipos:

\begin{itemize}
  \item \archivo{models.py}: registros, participantes, instituciones, equipos y miembros.
  \item \archivo{forms.py}: validaciones de registro de colegio y universidad.
  \item \archivo{views.py}: seleccion de tipo de registro, formularios, exito y confirmacion por token.
  \item \archivo{services.py}: correos de registro/asistencia, confirmacion y sincronizacion hacia \clase{Team}.
  \item \archivo{admin.py}: administracion de registros, equipos e instituciones; incluye acciones de solicitud de asistencia y sincronizacion.
  \item \archivo{templates/participants/}: pantallas publicas de registro y confirmacion.
\end{itemize}

Esta app conecta el mundo publico con el mundo competitivo. Por eso la regla de sincronizacion no queda solo en una vista: \funcion{confirm_attendance} delega creacion o actualizacion de equipos y puede activar inicializacion de competencia si corresponde.

\subsection{apps.tournament}

\archivo{apps.tournament} es el nucleo de competencia:

\begin{itemize}
  \item \archivo{models.py}: ediciones, reglas, fases, competencias, grupos, batallas, entradas, estados e historial.
  \item \archivo{formats.py}: perfiles automaticos y distribuciones por cantidad de equipos.
  \item \archivo{planner.py}: planificacion de divisiones y asignaciones.
  \item \archivo{recommendations.py}: recomendaciones de siguiente fase, repechaje y tercer lugar.
  \item \archivo{services.py}: inicializacion, grupos, batallas, repechajes, fases manuales, sincronizacion y podio.
  \item \archivo{views.py}: vista publica del torneo y panel interno.
  \item \archivo{admin.py}: administracion de modelos de torneo.
  \item \archivo{tests.py}: pruebas de distribuciones manuales.
  \item \archivo{management/commands/}: comandos \comando{bootstrap_tournament} y \comando{reset_tournament_demo}.
\end{itemize}

La concentracion de servicios en esta app facilita ubicar la logica competitiva. El riesgo asociado es que \archivo{services.py} tiene alta responsabilidad y debe modificarse con pruebas y revision de efectos colaterales.

\begin{advertencia}
El comando \comando{reset_tournament_demo} fue identificado como destructivo por la auditoria porque borra datos operativos antes de sembrar una simulacion. No debe tratarse como comando de mantenimiento rutinario ni ejecutarse en ambientes reales.
\end{advertencia}

\subsection{apps.invitations}

\archivo{apps.invitations} implementa comunicaciones:

\begin{itemize}
  \item \archivo{models.py}: plantillas, destinatarios, lotes y logs.
  \item \archivo{forms.py}: formularios de plantilla, destinatarios y lotes.
  \item \archivo{services.py}: marcadores DOCX, render, conversion PDF, XLSX, correo, logs y asistencia.
  \item \archivo{views.py}: dashboard, plantillas, previsualizacion, invitaciones, confirmaciones e historial.
  \item \archivo{admin.py}: administracion de plantillas, destinatarios, lotes y logs.
  \item \archivo{tests.py}: pruebas de conversion PDF y correo.
  \item \archivo{management/commands/}: comandos \comando{send_invitations} y \comando{send_attendance_requests}.
\end{itemize}

Esta separacion evita que el envio de comunicaciones quede disperso en participantes o torneo. Tambien permite implementar controles de simulacion, destinatario de prueba y logs por destinatario sin contaminar la logica competitiva.

\section{Templates y recursos estaticos}
\label{sec:estructura-templates-static}

Los templates se ubican dentro de cada aplicacion, siguiendo el dominio al que pertenecen:

\begin{longtable}{p{4cm}p{9cm}}
\toprule
Ruta & Uso \\
\midrule
\archivo{apps/core/templates/core/} & Base publica, home, reglas y patrocinadores. \\
\archivo{apps/participants/templates/participants/} & Registro, confirmacion de asistencia, errores y exito de registro. \\
\archivo{apps/tournament/templates/tournament/} & Vista publica del torneo, login interno, dashboard, division, fase, detalle y asistente manual. \\
\archivo{apps/invitations/templates/invitations/} & Dashboard de comunicaciones, plantillas, previsualizacion, invitaciones, confirmaciones e historial. \\
\bottomrule
\end{longtable}

Los recursos compartidos estan en \archivo{static/}:

\begin{itemize}
  \item \archivo{static/css/app.css}: estilos globales publicos e internos.
  \item \archivo{static/js/app.js}: comportamiento publico.
  \item \archivo{static/js/control.js}: panel interno, AJAX y modales.
  \item \archivo{static/docs/ExplotaGlobos.pdf}: documento publico.
  \item \archivo{static/sponsors/}: logos utilizados en vistas publicas.
\end{itemize}

Esta organizacion separa templates por dominio y assets por tipo. El mantenimiento del frontend debe revisar ambas zonas: el template que emite atributos \variable{data-*} y el JavaScript que los consume.

\section{Documentacion}
\label{sec:estructura-docs}

\archivo{docs/} contiene documentacion y soportes del proyecto:

\begin{itemize}
  \item \archivo{docs/auditoria_tecnica/}: auditoria tecnica validada.
  \item \archivo{docs/manual_tecnico/}: arquitectura y redaccion del presente manual.
  \item \archivo{docs/correos/}: insumos y ejemplos para comunicaciones.
  \item \archivo{docs/screenshots/}: capturas referenciadas por documentacion previa.
\end{itemize}

Los listados reales de comunicaciones y archivos sensibles dentro de \archivo{docs/correos/} estan protegidos por \archivo{.gitignore}. Esta regla evita versionar destinatarios reales o datos personales por accidente.

\section{Flujo entre aplicaciones}
\label{sec:estructura-flujo-apps}

El flujo entre aplicaciones puede entenderse en cuatro pasos:

\begin{enumerate}
  \item \archivo{apps.core} presenta informacion publica y direcciona al registro.
  \item \archivo{apps.participants} crea registros, confirma asistencia y transforma registros confirmados en equipos.
  \item \archivo{apps.tournament} consume equipos para construir competencia, fases, batallas y estados.
  \item \archivo{apps.invitations} se apoya en registros y confirmaciones para comunicaciones y trazabilidad de envios.
\end{enumerate}

\begin{figure}[H]
  \centering
  \fbox{\parbox{0.88\textwidth}{\centering Marcador de diagrama: flujo entre aplicaciones desde registro hasta competencia y comunicaciones. Fuente prevista: DIA-03 y DIA-07 en \archivo{planificacion/06_inventario_diagramas.md}.}}
  \caption{Marcador para diagrama de flujo entre aplicaciones del proyecto.}
  \label{fig:estructura-flujo-aplicaciones}
\end{figure}

\section{Buenas practicas observadas}
\label{sec:estructura-buenas-practicas}

La estructura aplica varias practicas utiles para mantenimiento:

\begin{itemize}
  \item Separacion por dominio mediante apps locales.
  \item Uso de servicios para logica que no debe vivir directamente en vistas.
  \item Plantillas ubicadas dentro de la app que las usa.
  \item Configuracion externa mediante variables de entorno.
  \item Exclusiones de \archivo{.gitignore} para secretos, bases locales, media, caches y listados reales.
  \item Comandos de gestion ubicados dentro de la app responsable.
  \item Tests colocados junto a los dominios que validan.
\end{itemize}

\section{Cierre del capitulo}
\label{sec:estructura-cierre}

La estructura del proyecto favorece una lectura por dominio: publico, participantes, torneo, comunicaciones y comun. Esta organizacion hace que el mantenimiento sea mas predecible porque permite rastrear cada flujo desde la ruta hasta la vista, servicio, modelo, template y asset relacionado. El punto que requiere mayor disciplina es el nucleo de torneo, donde la cantidad de responsabilidades de \archivo{apps/tournament/services.py} hace indispensable documentar contratos y ampliar pruebas antes de cambios profundos.
